\subsection{Class-preserving weight-five shadows and the odd-base \texorpdfstring{$A_6$}{A6} comparison}

\begin{statement}{D.38}{Proved}{Exact constraints on the first mixed tests above the characteristic five-group}
Every automorphism \(\beta_Q\) of \(P_1\) constructed in D.37 is
class-preserving.  There is an explicit conjugator for every element,
given in \textup{(D.38.3)} below.  The nonzero parameters nevertheless
retain their nonzero ordinary outer classes by the exact centralizer
calculation of D.37.

In addition, let
\[
 1\longrightarrow A_6^\Omega\longrightarrow E
       \xrightarrow{\pi}D\longrightarrow1
 \tag{D.38.1}
\]
be a finite chief extension with \(D\) of odd order.  Fix an
epimorphism \(q_E:\widehat F\twoheadrightarrow E\), write
\(X_i=q_E(x_i)\), \(q=\pi q_E\), \(d_i=q(x_i)\), and let
\(q_1:\widehat F\twoheadrightarrow D\) be any epimorphism with
\(q_1(x_i)=h_i d_i^\lambda h_i^{-1}\), for specified \(h_i\in D\)
and \(\lambda\in\widehat{\mathbf Z}^{\times}\).
For any lifts \(H_i\) of \(h_i\), there is an epimorphism
\(f:\widehat F\twoheadrightarrow E\) satisfying
\[
 \pi f=q_1,\qquad
 f(x_i)=a_iH_iX_i^\lambda H_i^{-1}a_i^{-1}
                  \quad(a_i\in A_6^\Omega).
 \tag{D.38.2}
\]
The number of these literal epimorphisms is the same as in the
original case \((q_1,\lambda,h_i)=(q,1,1)\).
The comparison also holds in the actual coupled pullback with
any finite odd-order group over \(D\), retaining its prescribed
peripheral epimorphism and conjugators.
In particular every \(\beta_Q\) passes every such chief step
over its fixed characteristic target \(P_1\).
\end{statement}

\begin{proof}
We first prove the elementwise conjugacy assertion in the actual
Magnus group, with the notation and multiplication of D.37.
Let \(g\in P_1\) have abelian coordinates \(r,s\) modulo five,
and write \(Q=AX^2+BXY+CY^2\).
If \(s=0\), then \(\beta_Q(g)=g\) by D.37.16.
Suppose \(s\ne0\), and put
\[
 \begin{split}
 d&=-sY(X+Y)Q,\\
 t&=d(s,-r)/s\in\mathbf F_5,\\
 p&=\frac{tsX^4-d}{rX+sY}\in\mathbf F_5[X,Y]_3,\\
 u_g&=x^{5t}c^p.
 \end{split}
 \tag{D.38.3}
\]
The quotient displayed here is an exact polynomial division.
Its homogeneous degree-four numerator vanishes at \((s,-r)\):
the first term evaluates to \(ts^5=ts=d(s,-r)\).
The kernel of substitution \(Y=-rX/s\) is the principal ideal
\((rX+sY)\), and the coefficient \(s\) is a unit.  This proves
existence and uniqueness of the cubic \(p\).

The following two commutators are exact in \(P_1\):
\[
 [x^5,g]=c^{sX^4},\qquad
 [c^p,g]=c^{-(rX+sY)p}.
 \tag{D.38.4}
\]
For the first, \([x,g]\) belongs to the polynomial derived module
and its constant coefficient is \(s\).  The product formula for
\([x^5,g]\) multiplies this coefficient polynomial by
\(1+(1+X)+\cdots+(1+X)^4=X^4\) in characteristic five.
Every positive-degree term then disappears at the cutoff \(I^5\).
For the second, conjugation by \(g\) on that module multiplies
by its diagonal \(a_g\); the linear part of \(a_g-1\) is
\(rX+sY\), and the terms of degree at least two kill the cubic
\(p\) at the same cutoff.  Both commutators in
\textup{(D.38.4)} are central.  Consequently
\[
 [u_g,g]=c^{tsX^4-(rX+sY)p}=c^d,
 \qquad u_ggu_g^{-1}=g c^d=\beta_Q(g).
 \tag{D.38.5}
\]
This proof includes all higher abelian coordinates and all fifteen
derived coordinates of \(g\); none was discarded from the group.
The term \(x^{5t}\) is essential.  For example, for
\(Q=X^2\), \((r,s)=(1,2)\), the value \(d(2,-1)\) is nonzero,
so a conjugator from the cubic derived module alone cannot give
the displayed central change.

The program \texttt{a\_weight5\_class\_preserving.py} evaluates
\textup{(D.38.3)--(D.38.5)} with actual Magnus multiplication
for all 125 parameters and all 25 abelian residue pairs, retaining
higher abelian and derived coordinates.  All 3,125 checks pass;
1,200 of them have nonzero displayed \(x^5\) coefficient.

We next give the complete odd-base chief comparison.  Set \(S=A_6\).
The actual automorphism enumeration of D.31 supplies
\(Z(S)=1\), simplicity, and
\(\operatorname{Out}(S)=\mathbf F_2^2\).
Let \(\rho:D\to\operatorname{Sym}(\Omega)\) be the transitive
factor action of \textup{(D.38.1)}.  Its outer action is
\((b_d,\rho(d))\), where
\(b_d\in(\mathbf F_2^2)^\Omega\) and
\(b_{de}=b_d+\rho(d)b_e\).
The explicit vector
\[
 B=|D|^{-1}\sum_{e\in D}b_e,
       \qquad b_d=B-\rho(d)B
 \tag{D.38.6}
\]
removes this twisting by componentwise actual automorphisms of
\(S\) having the prescribed outer coordinates.
The centreless pullback
\[
 E\longrightarrow
 \operatorname{Aut}(S^\Omega)
       \times_{\operatorname{Out}(S^\Omega)}D,
 \qquad e\longmapsto(\operatorname{Ad}_e,\pi(e))
 \tag{D.38.7}
\]
is an isomorphism: its kernel is \(Z(S^\Omega)=1\), and each
required automorphism differs from a chosen lift by a unique
inner automorphism.  After \textup{(D.38.6)}, the unique
coefficient \(a_e\) in
\(\tau_B^{-1}\operatorname{Ad}_e\tau_B
 =\operatorname{Inn}(a_e)\rho(\pi(e))\)
gives the actual isomorphism
\(e\mapsto(a_e,\pi(e))\) with \(S^\Omega\rtimes D\).
It transports the source and all specified conjugators.

Choose a factor stabilizer \(D_*\leq D\).
For every \(D_*\leq C\leq D\), use the two actual subgroups
\(H_C=q^{-1}(C)\), \(H'_C=q_1^{-1}(C)\), with respective
kernels \(K=\ker q\), \(K'=\ker q_1\).
The exact boundary bijection is
\[
 C d\langle q_1(x_i)\rangle
       \longmapsto C d h_i\langle d_i\rangle.
 \tag{D.38.8}
\]
It preserves each cycle length and every intermediate boundary
degree, because \(q_1(x_i)h_i=h_i d_i^\lambda\).
The induced cycle-product class is the \(\lambda\)th power
of the old class.  Kernel components of \(H_i\) only conjugate
that cycle product in \(S\).
For a boundary \(\delta\) of \(H_C\), let \(C_\gamma\)
be its specified classes upstairs and \(r_\gamma\) their
relative boundary degrees.  The exact allowable set is
\[
 T_{C,\delta}=
    \bigcap_{\gamma\mapsto\delta}
       \{s\in S:s^{r_\gamma}\in C_\gamma\}.
 \tag{D.38.9}
\]
The set on the new cover is \(T_{C,\delta}^\lambda\), with
the labels matched by \textup{(D.38.8)}; this equality retains
all roots when a boundary degree is not invertible on \(S\).

All unit-power class permutations of \(A_6\) are realized by
one actual automorphism \(\tau_\lambda\): use the identity
for \(\lambda\equiv1,4\pmod5\), and conjugation by \((01)\)
in \(S_6\) for \(\lambda\equiv2,3\pmod5\).
The only moved classes are the two classes of five-cycles;
the other cycle types each form one \(A_6\)-class.
Thus \(\tau_\lambda(T)=T^\lambda\) for every conjugacy-invariant
subset \(T\subset S\).
The program \texttt{a\_weight5\_a6\_odd.py} verifies this for
all sixteen units modulo \(\exp(A_6)=60\), on all 360 actual
permutations, and locates the literal odd conjugation in the
certified D.31 automorphism table.

The two covers have equal genus and matched oriented boundaries.
The adjacent Hurwitz maps and inverses displayed in D.29.13
therefore give an actual isomorphism \(u_C:H'_C\to H_C\)
with \textup{(D.38.8)} as its boundary correspondence.
The map
\[
                 \psi\longmapsto\tau_\lambda\psi u_C
 \tag{D.38.10}
\]
is a bijection on epimorphisms to \(S\) satisfying the sets
\textup{(D.38.9)}.  Denote the common count by \(N_C\).

Every epimorphism \(\varphi:H_{D_*}\to S\) maps \(K\)
onto \(S\): its normal image is either \(S\) or trivial,
and the latter would give an even-order quotient of the
odd-order group \(D_*\).  If
\(N=\ker(\varphi|_K)\), an overgroup \(H_C\) normalizes
\(N\) precisely when \(\varphi\) extends to \(H_C\).
Indeed its conjugation action on \(K/N\simeq S\) has outer
part factoring through the odd-order group \(C\), hence
trivial in \(\mathbf F_2^2\).
The inverse of \(S\xrightarrow{\operatorname{Inn}}
\operatorname{Inn}(S)\) supplies the extension.  It agrees
with \(\varphi\) and is unique by centrelessness and
surjectivity on \(K\).
Conversely each epimorphism from \(H_C\) restricts onto
\(S\) on \(K\), so every extension is counted exactly once.
Repeating the argument with \(K'\) proves that the numbers
with minimal kernel normalizer agree, by the explicit
finite subgroup-poset formula
\[
 e_{D_*}=\sum_{C\geq D_*}\mu(D_*,C)N_C.
 \tag{D.38.11}
\]
The original generating source makes this number positive.
The nonabelian Shapiro map of D.28.13 identifies the coordinate
kernels with the conjugates of \(N\); minimal normalizer makes
them pairwise distinct.  Surjective maps to a simple group
with distinct kernels have full product image: a proper
subdirect product would identify two factors by an
automorphism and hence identify their kernels.  This proves
generation of the new full kernel.
The literal gauge normalization of D.29.16, with \(S=A_6\),
has exactly \(360^{|\Omega|-1}\) choices, and retains the
specified base conjugators.  The exact literal count is
therefore
\[
       360^{|\Omega|-1}
           \sum_{C\geq D_*}\mu(D_*,C)N_C,
 \tag{D.38.12}
\]
which is unchanged under \textup{(D.38.10)}.

For a coupled odd-order group \(L\to D\), pair this lift
with its prescribed new map to \(L\), using the images of
its conjugators as the \(h_i\).  The paired conjugators
are actual elements of the fibre product.
The possible missing quotient of \(S^\Omega\) would be a
quotient of \(\ker(L\to D)\), of odd order.  Every nontrivial
quotient of \(S^\Omega\) contains an \(S\)-factor and has
even order, so no such quotient occurs.  The paired map
is therefore onto.  When \(L\) is a five-group, the
projection to \(L\) is the exact maximal pro-five quotient,
since the perfect kernel has trivial image in every
finite five-group.
\end{proof}

The single-kernel comparison in D.33 concerns odd-order
automorphisms over arbitrary finite bases, while D.34 and
D.35 use cyclic or two-group outer images of a centreless
kernel.  The comparison in D.38 instead handles every
transitive \(A_6\)-factor permutation action over an
odd-order base, arbitrary prescribed peripheral base
epimorphisms, and all unit multipliers.  Its maps
\textup{(D.38.6)} and \textup{(D.38.11)} explicitly use
oddness of the base.  D.36 has an even base and its
ordinary obstruction therefore remains in the further
mixed layers that D.38 does not decide.
